% DRAFT (SD) for the SD regions of devopsSre.tex (rows 1, 2, 3). Labels follow BD's skeleton (bd-drafts/devopsSre-skeleton.tex).
% Paths assume the manual repo layout: config/monitoring/... (copies of infra/monitoring) and evidence/NNN-....
% Before moving into the manual repo:
%   - copy infra/monitoring/{compose.yaml,prometheus/prometheus.yml,blackbox/blackbox.yml,Caddyfile.metrics-block.txt} to config/monitoring/
%   - copy evidence 110-114 (and BD's 140, 141, 143, 147 are cited) into evidence/
%   - copy sd-drafts/TikZ/devopsSreDeployment.tex to TikZ/devopsSre/
% Points per the course rubric (sec. 5.5): row 1 = 15, row 2 = 15, row 3 = 10.

% ===== SD region: rows 1, 2, 3 =====
% Region: sreDeployment (single owner; edit only your own region)
\section{Deployment and capacity (15 points)}\label{Section:sreDeployment}
\textbf{Owner:} SD. \textbf{Status:} deployed November 10, 2026; restart test November 11, 2026 \StatusPass.

\paragraph{Baseline and capacity.} Before anything was added, the server was read without changes (\ListRef{110}):
an EC2 \texttt{t3.micro} with 913~MiB of RAM, of which only 200~MiB was available and 425~MiB was already in swap, while
Gitea, Jenkins, Caddy and the Actions runner were running. Four more containers would not have fit. BD decided to resize
the same instance to \texttt{t3.small} (same disk and Elastic IP); afterwards the host had 1,909~MiB with 1,064~MiB
available and no swap in use (\ListRef{140}). The root disk had 8.4~GB free, enough for the images (about 1.4~GB) and
Prometheus' capped storage.

\paragraph{What runs.} Four separate containers, all from images published by the projects themselves; nothing is built
locally and no Dockerfile is written. Tags are pinned to exact releases (never \texttt{latest}) and the pulled image
digests are recorded (\ListRef{111} for the choice, \ListRef{112} for the digests on the server).

\begin{tabularx}{\linewidth}{L{2.9cm}L{4.6cm}Y}
\caption{Monitoring containers, versions, and limits\label{Table:SreContainers}}\\
\toprule \textbf{Container} & \textbf{Image} & \textbf{Role, limit, storage}\\ \midrule
\endfirsthead
\toprule \textbf{Container} & \textbf{Image} & \textbf{Role, limit, storage}\\ \midrule
\endhead \bottomrule \endlastfoot
\texttt{prometheus} & \texttt{prom/prometheus:}\newline\texttt{v3.16.0} & scrapes every 30\,s and stores metrics; 256~MiB; volume \texttt{monitoring\_prometheus\_data}, kept 7 days or 1~GB, whichever comes first\\
\texttt{grafana} & \texttt{grafana/grafana:}\newline\texttt{13.3.1} & dashboard and alerts; 384~MiB; volume \texttt{monitoring\_grafana\_data} (users, data source, dashboard, alert rules)\\
\texttt{node-exporter} & \texttt{quay.io/prometheus/}\newline\texttt{node-exporter:v1.12.1} & whole-host CPU, memory, root filesystem; 64~MiB; host \texttt{/} mounted read-only\\
\texttt{blackbox-}\newline\texttt{exporter} & \texttt{quay.io/prometheus/}\newline\texttt{blackbox-exporter:v0.29.1} & HTTPS probes of the Gitea and Jenkins login pages; 64~MiB; no storage\\
\end{tabularx}

The memory limits make sure that, if memory runs out, a monitoring container is stopped rather than Gitea or Jenkins.
Grafana's limit started at 256~MiB, but Grafana used 216~MiB while idle, so it was raised to 384~MiB on the first day
(recorded in \ListRef{112}). Retention of 7 days follows the course's lab suggestion; the 1~GB cap protects the root
disk that Jenkins builds also use.

\paragraph{Configuration in Gitea.} The non-secret files are in the \texttt{infra} repository, folder
\href{\GiteaURL/group1-startup/infra/src/branch/main/monitoring}{\texttt{monitoring}}, added through
\href{\GiteaURL/group1-startup/infra/issues/1}{issue \#1} and
\href{\GiteaURL/group1-startup/infra/pulls/2}{pull request \#2}. BD checked the Caddy block against the live
Caddyfile and the memory limits against the running containers, and merged it on November 11, 2026, 19:42~EST, as
merge commit \href{\GiteaURL/group1-startup/infra/commit/c03fb1b10e003554374404b7759cefa408214314}{\texttt{c03fb1b}}
(her review was given in the team chat; Gitea records the merge, not a formal approval). The same files run in \texttt{\textasciitilde/monitoring} on
the server. No Compose service publishes a host port, none is privileged, and none mounts the Docker socket.

\begin{longlisting}
\inputminted[fontsize=\scriptsize]{yaml}{config/monitoring/compose.yaml}
\caption{\texttt{compose.yaml} of the monitoring stack (\texttt{infra/monitoring}): four containers, no published ports, memory limits, named volumes, Grafana also on the existing \texttt{gitea\_web} network.}
\label{Listing:ESreCompose}
\end{longlisting}

\paragraph{Start, update, recover.} From \texttt{\textasciitilde/monitoring} on the server:
\begin{minted}{bash}
cp .env.example .env                        # image tags only
docker compose config --quiet && docker compose pull && docker compose up -d
docker compose ps                           # four containers Up
\end{minted}
To update, change a tag in \texttt{.env}, run \texttt{docker compose pull \&\& docker compose up -d}, and record the new
digests. To stop and start, use \texttt{docker compose stop} and \texttt{docker compose start}. \textbf{Never}
\texttt{docker compose down -v}: it deletes both volumes, that is, the metric history and all of Grafana's state.
The first Grafana administrator password was a random value generated on the server into
\texttt{secrets/grafana\_admin\_password} (folder mode 700, never printed, not in Git), so the default
\texttt{admin}/\texttt{admin} was never valid on the public login page; BD then set her own password with
\texttt{grafana cli admin reset-admin-password} (Section~\ref{Section:sreAccess}).

\paragraph{Deployment diagram.} \FigRef{SreDeployment} adapts the course's Figure~5.1 to this server.

\begin{figure}[htbp]\centering
\resizebox{\linewidth}{!}{\input{TikZ/devopsSre/devopsSreDeployment.tex}}
\caption{Deployment of the Group 1 server with monitoring, adapted from the course manual's Figure~5.1. Blue: the four new
monitoring containers. Arrows start at the side that opens the connection. Gitea and Jenkins are probed only through their
public HTTPS names; Prometheus, the exporters, and Grafana's port 3000 are reachable only on private Docker networks.
Not drawn: the Gitea Actions runner and the Gitea--Jenkins webhook (Chapter~\ref{Chapter:Jenkins}).}
\label{Fig:ESreDeployment}
\end{figure}

\paragraph{Restart persistence.} Right after BD's controlled Jenkins outage and after the SLI window had been read
(Section~\ref{Section:sreSlo}), Prometheus and then Grafana were restarted with \texttt{docker restart}, on November 11,
2026, 8:38--8:40~PM New York time (\ListRef{115}). Both kept the same container IDs and the same named volumes. Prometheus'
oldest stored sample (03:15:08~UTC on November 11, the first scrape) and its 720 probe samples per URL from the last six
hours were still there; Prometheus was ready again after 3~seconds, and one 30-second Gitea sample was missed (719 instead
of 720). Grafana was healthy after 14~seconds, and its data source, both alert rules, the contact point
\texttt{team-alerts}, both user accounts, and the dashboard (created at 8:25~PM, before the restart) were all still
present. A first reading reported zero dashboards: the script had queried Grafana's legacy dashboard table, while
Grafana~13 keeps dashboards in its newer \texttt{resource} table; the corrected check is in the same listing. This shows
that the volumes persist; it is not a backup and restore. The same test is BD's verification in
Section~\ref{Section:sreIncident}.

\begin{longlisting}
\inputminted[fontsize=\scriptsize]{text}{evidence/110-server-baseline-before-monitoring-sd-2026-11-10.txt}
\caption{Baseline before monitoring: instance type, memory and swap, disk, containers, listening ports, November 10, 2026 (real output).}
\label{Listing:E110}
\end{longlisting}

\begin{longlisting}
\inputminted[fontsize=\scriptsize]{text}{evidence/111-image-tags-digests-amd64-and-local-config-validation-sd-2026-11-10.txt}
\caption{Chosen tags with index and linux/amd64 digests, and local validation of the configuration, November 10, 2026 (real output).}
\label{Listing:E111}
\end{longlisting}

\begin{longlisting}
\inputminted[fontsize=\scriptsize]{text}{evidence/112-monitoring-stack-deployed-digests-targets-probes-host-values-sd-2026-11-10.txt}
\caption{Deployment: pull and start, container settings, pulled digests, Prometheus targets, first probe results, host values, Grafana limit change, November 10, 2026 (real output).}
\label{Listing:E112}
\end{longlisting}

\begin{longlisting}
\inputminted[fontsize=\scriptsize]{text}{evidence/115-prometheus-grafana-restart-persistence-test-sd-2026-11-11.txt}
\caption{Restart test of Prometheus and Grafana: state before and after, readiness times, collection gap, and the corrected dashboard check, November 11, 2026 (real output).}
\label{Listing:E115}
\end{longlisting}

% Region: sreMetrics (single owner; edit only your own region)
\section{Metrics and HTTPS probes (15 points)}\label{Section:sreMetrics}
\textbf{Owner:} SD. \textbf{Status:} all targets up and host values verified November 10, 2026 \StatusPass.

\paragraph{What is scraped.} Prometheus scrapes four jobs every 30~seconds, all by private service name on the
\texttt{monitoring} network (\ListRef{SrePrometheus}): itself (\texttt{prometheus:9090}), the Node Exporter
(\texttt{node-exporter:9100}), the Blackbox Exporter's own metrics (\texttt{blackbox-exporter:9115}), and the
\texttt{https} probe job. Grafana reads only from Prometheus; it does not contact Gitea, Jenkins, or the exporters.
Right after the start, all five targets were \texttt{up} with no scrape error (\ListRef{112}).

\begin{longlisting}
\inputminted[fontsize=\scriptsize]{yaml}{config/monitoring/prometheus/prometheus.yml}
\caption{\texttt{prometheus.yml}: four jobs, 30\,s interval, and the multi-target relabelling that sends each URL to the Blackbox Exporter.}
\label{Listing:ESrePrometheus}
\end{longlisting}

\paragraph{HTTPS probes and what ``success'' means.} The \texttt{https} job asks the Blackbox Exporter to fetch
\texttt{https://git.ssw590-g1.tech/user/login} and \texttt{https://ci.ssw590-g1.tech/login}. The relabelling turns each
URL into the \texttt{target} parameter and the \texttt{instance} label, and sends the request to
\texttt{blackbox-exporter:9115}. A probe counts as successful (\texttt{probe\_success} = 1) only if the public DNS name
resolves over IPv4, the TLS certificate is valid for that name (verification is on), the answer is HTTPS, and the final
status after redirects is 2xx, all within 10~seconds (\ListRef{SreBlackbox}). The probe uses the same public names as a
user: it leaves through the Elastic IP and comes back in through the security group and Caddy (\FigRef{SreDeployment}).
The first probes returned success 1, HTTP 200, in 0.014\,s (Gitea) and 0.026\,s (Jenkins) (\ListRef{112}).

\begin{longlisting}
\inputminted[fontsize=\scriptsize]{yaml}{config/monitoring/blackbox/blackbox.yml}
\caption{\texttt{blackbox.yml}: module \texttt{http\_2xx} with TLS verification on and plain HTTP counted as failure.}
\label{Listing:ESreBlackbox}
\end{longlisting}

\paragraph{Metric inventory.}
\begin{tabularx}{\linewidth}{L{4.2cm}L{1.3cm}Y}
\caption{Metrics used in this chapter\label{Table:SreMetrics}}\\
\toprule \textbf{Metric} & \textbf{Type} & \textbf{Meaning and limitation}\\ \midrule
\endfirsthead
\toprule \textbf{Metric} & \textbf{Type} & \textbf{Meaning and limitation}\\ \midrule
\endhead \bottomrule \endlastfoot
\texttt{node\_cpu\_seconds\_total} & counter & CPU seconds per core and mode since boot; only meaningful through \texttt{rate()}; whole host, not per container\\
\texttt{node\_memory\_MemAvailable\_bytes} & gauge & the kernel's estimate of memory that programs can still get (not ``free'' memory)\\
\texttt{node\_filesystem\_avail\_bytes}, \texttt{node\_filesystem\_size\_bytes} & gauge & bytes usable by normal users, and size, per filesystem; the root filesystem must be selected by its labels\\
\texttt{up} & gauge & 1 if Prometheus' last scrape of a target worked; says nothing about the probed service\\
\texttt{probe\_success} & gauge & 1 if the whole HTTPS check passed; a login page only, not a Git push or a build\\
\texttt{probe\_duration\_seconds} & gauge & time of the whole probe, including DNS, TLS, and redirects\\
\end{tabularx}

\paragraph{Host values checked against the operating system.} At 03:16~UTC on November 11 (evening of November 10,
Eastern time), the Prometheus values were compared with \texttt{free -b} and \texttt{df -B1 /} run on the host
(\ListRef{112}). The root filesystem series has \texttt{mountpoint="/"}, \texttt{device="/dev/nvme0n1p1"}, \texttt{fstype="xfs"};
overlay and tmpfs mounts are excluded, never averaged in.

\begin{tabularx}{\linewidth}{L{3.6cm}L{3.4cm}L{3.4cm}Y}
\caption{Prometheus compared with the host\label{Table:SreHostCheck}}\\
\toprule \textbf{Value} & \textbf{Prometheus} & \textbf{Host command} & \textbf{Result}\\ \midrule
\endfirsthead
\toprule \textbf{Value} & \textbf{Prometheus} & \textbf{Host command} & \textbf{Result}\\ \midrule
\endhead \bottomrule \endlastfoot
Total memory & 2,002,350,080 & \texttt{free -b}: 2,002,350,080 & equal\\
Available memory & 823,209,984 & \texttt{free -b}: 804,671,488 & 2\,\% apart; sampled about 10\,s apart\\
Root size & 21,395,124,224 & \texttt{df -B1}: 21,395,124,224 & equal\\
Root available & 6,994,649,088 & \texttt{df -B1}: 6,994,644,992 & one 4~KiB block apart\\
CPUs & 2 idle series & 2 vCPUs (\texttt{t3.small}) & equal\\
\end{tabularx}

Root free space was therefore 32.7\,\% (68\,\% used, up from 58\,\% because of the new images), well above the 15\,\%
alert threshold of Section~\ref{Section:sreAlerts}.

\paragraph{Limitations.} A 30-second scrape cannot see events shorter than about 30~seconds. The probes run on the same
server they check, so a total server outage stops the probes and the alerts too; an independent outside check is proposed
in Section~\ref{Section:sreIncident}. A probe from AWS does not prove that every user network can reach the services.
Node Exporter reports the whole host; per-container usage (cAdvisor) is outside the required scope.

% Region: sreAccess (single owner; edit only your own region)
\section{Access control and exposure (10 points)}\label{Section:sreAccess}
\textbf{Owner:} SD (Caddy block BD, accounts BD). \textbf{Status:} tested from a team and a non-team network \StatusPass.

\paragraph{Login URLs.} All three services use HTTPS with Let's Encrypt certificates; no password is given here
(see the Hosts and Passwords chapter for where each is kept):
\begin{itemize}[nosep]
  \item Gitea: \url{https://git.ssw590-g1.tech/user/login}
  \item Jenkins: \url{https://ci.ssw590-g1.tech/login}
  \item Grafana: \url{https://metrics.ssw590-g1.tech/login} (certificate issued November 11, 2026, \ListRef{113})
\end{itemize}

\paragraph{How Grafana is reached.} The DNS name \texttt{metrics.ssw590-g1.tech} points to the Elastic IP (\ListRef{141}).
Caddy, the existing TLS proxy, is the only container with public ports (443 and 80). It forwards to \texttt{grafana:3000}
on the private \texttt{gitea\_web} network. The block follows the Jenkins rule: addresses on the team list
(\texttt{TEAM\_IPS}, kept only in the server's \texttt{.env}) reach all of Grafana; everyone else reaches only the login
page, its static files, and the health check; everything else gets \texttt{403 Forbidden}.

\begin{longlisting}
\inputminted[fontsize=\scriptsize]{text}{config/monitoring/Caddyfile.metrics-block.txt}
\caption{Caddy site block for Grafana, as applied on the server (\texttt{infra/monitoring/Caddyfile.metrics-block.txt}).}
\label{Listing:ESreCaddy}
\end{longlisting}

\paragraph{Authentication.} Grafana requires a sign-in: anonymous access and self sign-up are off (\ListRef{SreCompose}).
The administrator account belongs to BD; its first password was random and server-only, and BD replaced it (\FigRef{144}).
SD has her own named account with the Editor role, not an administrator (\ListRef{147}); her temporary password was handed
over privately, and SD replaced it with her own on November 11, 2026, about 20:15~EST, signed in from the campus network (on the team list). A grader account is not needed: the dashboard and
alert evidence is shown in Sections~\ref{Section:sreDashboard} and~\ref{Section:sreAlerts}.

\paragraph{Exposure tests.} The security group is unchanged for monitoring: only 443, 80, and 22 (team addresses, key
login only) are open. Ports 3000, 9090, 9100 and 9115 were closed from the internet (\ListRef{113}, \ListRef{114}).
The decisive test was from a network that is on no team list, a T-Mobile mobile connection (\ListRef{114}):

\begin{tabularx}{\linewidth}{L{5.6cm}L{2.0cm}Y}
\caption{Grafana from a non-team network, November 11, 2026\label{Table:SreOutside}}\\
\toprule \textbf{Request} & \textbf{Result} & \textbf{Meaning}\\ \midrule
\endfirsthead
\toprule \textbf{Request} & \textbf{Result} & \textbf{Meaning}\\ \midrule
\endhead \bottomrule \endlastfoot
\texttt{/login} and its 9 static files & 200 & the login page works for anyone, certificate verified\\
\texttt{/api/health} & 200 & health check only\\
\texttt{/}, \texttt{/explore}, \texttt{/d/...} & 403 & no dashboards or queries\\
\texttt{/api/datasources}, \texttt{/api/user}, \texttt{/api/search}, \texttt{/api/admin/users} & 403 & no API\\
\texttt{/metrics}, alert API & 403 & Grafana's own metrics and alerts hidden\\
\texttt{POST /login} as \texttt{admin}/\texttt{admin} & 401 & the default password never worked\\
TCP 22, 3000, 9090, 9100, 9115 & closed & SSH limited to team addresses; monitoring ports private\\
\end{tabularx}

\paragraph{Limitations.} The login form is public, so anyone can try passwords; Grafana's built-in protection against
repeated failed logins applies. The campus network's address is on the SSH and Caddy team lists from November 11, 2026
until teardown; it is shared by everyone on campus, so it narrows access less than a personal address. SSH still needs a
key, and Grafana still needs a password.

\begin{longlisting}
\inputminted[fontsize=\scriptsize]{text}{evidence/113-grafana-https-access-control-team-and-non-team-sd-2026-11-10.txt}
\caption{Grafana HTTPS: certificate, team view, simulated non-team view, closed ports, November 10, 2026 (real output).}
\label{Listing:E113}
\end{longlisting}

\begin{longlisting}
\inputminted[fontsize=\scriptsize]{text}{evidence/114-grafana-access-from-non-team-mobile-network-sd-2026-11-11.txt}
\caption{Grafana from a non-team mobile network (T-Mobile): only the login page is reachable, November 11, 2026 (real output).}
\label{Listing:E114}
\end{longlisting}
